%% shortname: Demand, Walras Law and WARP
\documentclass[11pt]{article}
\usepackage{microstyle}

\begin{document}

\lessonheader{3}{Chapter 2 (continued): demand, Walras' law, and WARP}{}


\section{Demand functions}

The consumer's choice depends on prices $p$ and wealth (or income) $w$. We let
$x(p,w)$ denote the \emph{demand correspondence}: the set of bundles the consumer
chooses from the budget set $B_{p,w}$. If the correspondence is single-valued --- that is, the
set consists of a single element --- then we have a standard demand function. In
the latter case
\[ x(p,w)=\bigl(x_1(p,w),\dots,x_L(p,w)\bigr), \]
where $x_i(p,w)$ is the individual demand for good $i$, $i=1,2,\dots,L$.

We start with some basic assumptions, to the effect that the consumer is rational.

\topic{Homogeneity of degree 0}
$x(p,w)$ is homogeneous of degree 0:
\[ x(\alpha p,\alpha w)=x(p,w)\qquad\text{for all }\alpha>0. \]
This is saying that only relative prices and real income matter (the budget set
is the same, $B_{\alpha p,\alpha w}=B_{p,w}$). It does not matter if
we convert to a different currency, or from dollars to cents.

\section{Walras' law and comparative statics}

\topic{Walras' law}
For any $p\gg 0$ and $w>0$,
\[ p\cdot x(p,w)=w, \]
i.e.\ the agent uses all his wealth (income).

We now assume that the demand correspondence is a demand function. We often assume
it is differentiable.

\topic{Comparative statics}
The parameters are $(p,w)$. What happens to the endogenous choices (i.e.\
consumption) when the exogenous parameters change? With differentiable demand
functions we want to examine
\[ \pd{x_i(p,w)}{p_j}\qquad\text{and}\qquad \pd{x_i(p,w)}{w}. \]

We now use the assumptions for comparative statics. Using homogeneity of degree 0,
\[ x(\alpha p,\alpha w)=x(p,w)\ \text{ for }\ \alpha>0
   \qquad\Longrightarrow\qquad
   \frac{d\,x_i(\alpha p,\alpha w)}{d\alpha}=0, \]
where the derivative on the left is the \emph{total} derivative with respect to
$\alpha$, i.e.\ the chain rule applied to all the arguments
$\alpha p_1,\dots,\alpha p_L,\alpha w$. Hence
\[ 0=\sum_{j=1}^{L}\pd{x_i(\alpha p,\alpha w)}{p_j}\,p_j
     +\pd{x_i(\alpha p,\alpha w)}{w}\,w . \]
Evaluate at $\alpha=1$:
\[ 0=\sum_{j=1}^{L}\pd{x_i(p,w)}{p_j}\,p_j
     +\pd{x_i(p,w)}{w}\,w . \]
When $x_i(p,w)>0$, multiply the whole expression by $x_i(p,w)$ and divide each term
inside by the same $x_i(p,w)$ --- the factor cancels, so the equality is untouched:
\[
\begin{aligned}
  0&=\sum_{j=1}^{L}\pd{x_i(p,w)}{p_j}\,p_j
      +\pd{x_i(p,w)}{w}\,w\\[2pt]
   &=x_i(p,w)\left[\sum_{j=1}^{L}\pd{x_i(p,w)}{p_j}\,\frac{p_j}{x_i(p,w)}
      +\pd{x_i(p,w)}{w}\,\frac{w}{x_i(p,w)}\right].
\end{aligned}
\]
The bracket is the sum of the demand elasticities of good $i$ with respect to every
price (own and cross) and to wealth, and since $x_i(p,w)>0$ it must itself be zero.

This is intuitive, given that the starting point is homogeneity of degree 0:
nothing changes if all prices and wealth change by the same percentage. We learn
that the elasticities cannot all have the same sign: unless they are all zero, some
are positive and some negative.

\section{The weak axiom of revealed preference (WARP)}

WARP is a minimal rationality assumption: the agent does not make
self-contradictory choices.

WARP is satisfied if, for any two pairs $(p,w)$ and $(p',w')$, it holds that
\[ p\cdot x(p',w')\le w
   \quad\text{and}\quad x(p',w')\ne x(p,w)
   \qquad\Longrightarrow\qquad
   p'\cdot x(p,w)>w' . \]
Tomorrow's bundle is affordable today, but I chose not to buy it: $x(p,w)$ is
revealed preferred to $x(p',w')$. The reason I do not buy $x(p,w)$ tomorrow is
simply that I cannot afford it.

\begin{figure}[ht]
\centering
\begin{tikzpicture}[scale=0.40]
  % ---------------- panel (a): WARP holds ----------------
  \begin{scope}
    \ecaxes{12.5}{8.5}
    \draw[budget] (0,8) -- (6.5,0);                 % B(p,w) steeper
    \draw[budget] (0,4.5) -- (11.5,0);              % B(p',w') flatter

    \node[ecnode,anchor=south west] at (2.5, 5.27) {$B_{p,w}$};
    \node[ecnode,anchor=south west] at (8.0, 1.8) {$B_{p',w'}$};

    \node[bundle] at (1.63, 6.0) {};                % x on B(p,w)
    \node[ecnode,anchor=south west] at (1.7, 6.1) {$x$};

    \node[bundle] at (3.83, 3.0) {};                % x' on B(p',w'), inside B(p,w)
    \node[ecnode,anchor=north east] at (3.8, 2.7) {$x'$};
    \node[ecnode,anchor=north] at (6.25,-0.6) {WARP holds};
  \end{scope}
  % ---------------- panel (b): WARP fails ---------------
  \begin{scope}[xshift=15cm]
    \ecaxes{12.5}{8.5}
    \draw[budget] (0,4.5) -- (11.5,0);              % B(p,w) flatter
    \draw[budget] (0,7.5) -- (6.5,0);               % B(p',w') steeper

    \node[ecnode,anchor=south west] at (8.0, 1.8) {$B_{p,w}$};
    \node[ecnode,anchor=south west] at (1.5, 6.0) {$B_{p',w'}$};

    \node[bundle] at (2.56, 3.5) {};                % x on B(p,w), inside B(p',w')
    \node[ecnode,anchor=north west] at (2.66, 3.0) {$x$};

    \node[bundle] at (4.77, 2.0) {};                % x' on B(p',w'), inside B(p,w)
    \node[ecnode,anchor=north east] at (4.6, 1.9) {$x'$};
    \node[ecnode,anchor=north] at (6.25,-0.6) {WARP fails};
  \end{scope}
  \end{tikzpicture}
\caption{The two cases for a pair of price--wealth situations $(p,w)$ and
$(p',w')$. \emph{Left:} $x'$ is affordable when $x$ is chosen under $(p,w)$, revealing
that $x$ is preferred to $x'$. Because $x$ is unaffordable under $(p',w')$, there is no
contradiction and WARP holds. \emph{Right:} $x'$ is affordable when $x$ is chosen
under $(p,w)$, and $x$ is affordable when $x'$ is chosen under $(p',w')$. This creates
a mutual revealed preference, which is exactly what WARP forbids.}
\end{figure}

\section{WARP and Slutsky wealth compensation}

WARP implies homogeneity of degree 0. If $(p',w')=(\alpha p,\alpha w)$, then $x'$
is affordable whenever $x$ was chosen, and vice versa. If $x'\ne x$, then $x$ is
revealed preferred to $x'$ and $x'$ is revealed preferred to $x$ --- a
self-contradictory choice, which violates WARP. Hence $x'=x$.

WARP has strong implications when combined with Walras' law.

\topic{A thought experiment}
Take a price change from $p$ to $p'$ and, at the same time, adjust income so that
\[ w'=p'\cdot x(p,w), \]
which is exactly \emph{Slutsky wealth compensation}: at the new prices the original
bundle is just affordable. The result (the \emph{compensated law of demand}): suppose $x(p,w)$ is homogeneous
of degree 0 and satisfies Walras' law. Then $x(p,w)$ satisfies WARP if and only if,
for any compensated price change from $(p,w)$ to $(p',w')$, with
$w'=p'\cdot x(p,w)$, it holds that
\[ (p'-p)\cdot\bigl(x(p',w')-x(p,w)\bigr)\le 0, \]
with strict inequality if $x(p,w)\ne x(p',w')$.

In particular, if only the price of good $i$ changes ($p_j'=p_j$ for $j\ne i$), the
inner product has a single term,
\[ (p_i'-p_i)\bigl(x_i(p',w')-x_i(p,w)\bigr)\le 0, \]
so the compensated demand for good $i$ moves in the opposite direction to its own
price. When several prices change at once only the inner product is signed, not each
good separately.

\begin{figure}[ht]
\centering
\begin{tikzpicture}[scale=0.52]
  \ecaxes{9.6}{6.6}
  \draw[budget] (0,6) -- (6,0);                   % B(p,w)
  \draw[budget] (0,3) -- (9,0);                   % B(p',w') : w' = p'.x(p,w)
  \draw[adm]    (4.5,1.5) -- (9,0);               % marked: the part with p.x(p',w') > w
  \node[ecnode,anchor=south west] at (1.40,4.70) {$B_{p,w}$};
  \node[ecnode,anchor=south west] at (0.60,2.85) {$B_{p',w'}$};
  \node[bundle] at (4.5,1.5) {};                  % x(p,w), at the crossing
  \node[ecnode,anchor=north east] at (4.35,1.30) {$x(p,w)$};
  \node[ecnode] at (7.00,2.20) {$p\cdot x'>w$};
  \draw[thin] (6.90,1.90) -- (6.50,0.86);
  \node[bundle] at (7.5,0.5) {};                  % x' = x(p',w'), on the marked part
  \node[ecnode,above right] at (7.62,0.60) {$x'$};
\end{tikzpicture}
\caption{Slutsky wealth compensation. The new budget line $B_{p',w'}$ is drawn so
that $w'=p'\cdot x(p,w)$: it passes through the original bundle $x(p,w)$, so the two
lines cross exactly there. Write $x'=x(p',w')$ for the bundle chosen at the new
prices. Since $x'$ must satisfy $p'\cdot x'=w'$ it lies on $B_{p',w'}$, and since
WARP requires $x'$ to be unaffordable at $(p,w)$ it can only lie on the thick part
of that line, where $p\cdot x'>w$.}
\end{figure}

\topic{Reading the picture}
Because $p'\cdot x(p,w)=w'$, the original bundle lies on $B_{p',w'}$, so the two
lines cross at $x(p,w)$. The new choice $x'=x(p',w')$ also lies on $B_{p',w'}$
(Walras' law). If $x'\ne x(p,w)$ were on the stretch inside $B_{p,w}$ (from the
crossing up to the vertical intercept), each bundle would be affordable when the
other was chosen, which WARP forbids. So $x'$ lies on the thick stretch, where
$p\cdot x'>w$. Here good 2 has become relatively more expensive (the new line is
flatter), and the thick stretch has less of good 2 than $x(p,w)$: demand for the good
whose price rose falls.

\begin{transcriptnotes}
  \item The original opens with a divider page reading only ``Chapter 2''; it is
        represented here by the subtitle.
  \item ``single-valuable'' $\to$ \emph{single-valued}; ``weath'' $\to$
        \emph{wealth}; ``connect to different currency'' $\to$ \emph{convert}.
  \item ``Walva's'' $\to$ \emph{Walras'}; ``homogenity'' $\to$
        \emph{homogeneity}; ``as verse versa'' $\to$ \emph{and vice versa}.
  \item ``the agent uses all his wealth on income'' $\to$ \emph{uses all his wealth
        (income)}.
  \item The step after evaluating at $\alpha=1$ is kept as in the original: multiply
        the whole expression by $x_i(p,w)$ and divide each term inside by the same
        $x_i(p,w)$. The two occurrences of the factor cancel, so the equality $=0$
        survives, and because $x_i(p,w)>0$ the bracket --- the sum of the
        elasticities --- must itself be zero.
  \item \textbf{Corrected:} ``all income elasticities sum to zero'' is too narrow.
        What homogeneity gives is that the price elasticities of the demand for
        good $i$ (own and cross) \emph{plus} its income (wealth) elasticity sum to
        zero.
  \item \textbf{Corrected:} the compensated-change condition was written with
        ``$=0$''. Since the next line says ``strict if $x(p,w)\ne x(p',w')$'', the
        intended relation is $(p'-p)\cdot(x(p',w')-x(p,w))\le 0$: the equality sign
        was a slip.
  \item \textbf{Corrected:} the line ``in the aggregate case'' stated the sign
        condition good by good, $(p_i'-p_i)(x_i'-x_i)\le0$. That does not follow
        when several prices change together: only the inner product over all goods
        is signed. The good-by-good statement holds when only $p_i$ changes, and it
        is written with that qualifier.
  \item \textbf{Corrected:} the statement of the compensated law of demand omitted
        homogeneity of degree 0, which the ``if and only if'' assumes (MWG
        Proposition 2.F.1); it is added. The sentence ``for the
        compensated demand, assume that $x(p,w)$ satisfies Walras' law'' is
        folded into that statement.
  \item \textbf{Corrected:} ``we learn that some elasticities are positive and some
        negative'' is stated too strongly (they may all be zero); what follows from
        a zero sum is that they cannot all share a strict sign.
  \item ``the set of bundles that are acceptable'' is written as the set of bundles
        \emph{chosen} from $B_{p,w}$; ``real prices'' is written \emph{relative
        prices}. The derivative in $\alpha$ is typeset $d/d\alpha$, since it is a
        total derivative.
  \item The last sketch was redrawn with $x(p,w)$ at the crossing of the two budget
        lines --- with Slutsky compensation the new line must pass through the
        original bundle --- and with the stretch of $B_{p',w'}$ on which
        $x'=x(p',w')$ is admissible, i.e.\ where $p\cdot x'>w$, drawn thick and
        labelled, and $x'$ shown on it.
  \item The summation indices are written inconsistently in the original
        (``$\sum_{i=1}^{L}\partial x_i(\alpha p,\alpha w)/\partial p_i\,p_j$'',
        ``$\partial p_2\,p_2$'' after setting $\alpha=1$, and
        ``$\sum_i \partial x_i/\partial p_i\cdot p_i/x_i$'' in the elasticity line, which
        would leave only the own-price term); the wealth derivative is written
        ``$\partial x_i(\alpha p,w)/\partial w$''. All are typeset with the sum over
        $j$ and the argument $(\alpha p,\alpha w)$.
  \item \textbf{Added:} the paragraph ``Reading the picture'' (why $x'$ must lie on the
        thick stretch, and that the stretch has less of good 2) is an explanation of
        the final sketch, not written in the original.
  \item The WARP figure now follows the corrected reading. In the holds case only
        $x'$ is affordable at $(p,w)$ while $x$ is unaffordable at $(p',w')$, so the
        antecedent of the axiom is met and its conclusion holds; in the failing case
        both bundles are affordable under both price--wealth pairs. The two panels
        use the same two lines with the labels swapped ($B_{p,w}$ the steeper line on
        the left, the flatter one on the right), so the two cases are read against the
        same geometry. The point labels are placed just clear of the lines. The panels are captioned ``WARP holds'' and ``WARP fails'' (the original labels only the second).
\end{transcriptnotes}

\end{document}
